\subsection{Independent removal attacks}
\label{sec:preserving_attack}
After both prediction-anchored models pass the initial quality screen, independent attackers run 500 clean AdamW updates on 4,000 fitting images from classes 400--799. The simulated adaptation used 2,000 different images from classes 0--199. Both image identities and class labels are disjoint between those support and attack pools. The reserved final classes are outside all 800 fitting classes; they are not claimed to be unseen during ImageNet pretraining or earlier project experiments.

The attacker receives the supplied marked model and a frozen copy of it as a prediction anchor, but neither the owner's detector and message nor the original clean weights. It trains with anchor coefficients $0,1,4,16$ and saves updates 25, 100 and 500. All 12 candidates per model are evaluated on 32 coupled validation seeds using the original preservation metrics in Subsection~\ref{sec:preservation}. The frozen rule selects the largest update count with at least 28/32 quality passes, then the smallest anchor coefficient. It chooses update 25 with coefficient 16 for the released model, and update 25 with coefficient 4 for each trained model. The selections are frozen before generation and detector scoring on 64 different final seeds.

\begin{table}[ht]\centering\small
\begin{tabular}{lrrrrr}\toprule
Starting model & Updates & Anchor & Preserved & Raw removed & Joint removed\\\midrule
Original & 25 & 16.0 & 60/64 & 0/64 & 0/64\\
Fixed & 25 & 4.0 & 56/64 & 0/64 & 0/64\\
Adapted & 25 & 4.0 & 57/64 & 0/64 & 0/64\\
\bottomrule\end{tabular}
\caption{Frozen quality-only-selected image attacks on 64 reserved final seeds and fitting-holdout classes. Each attack is paired with its own supplied marked checkpoint. The distinct selected anchor weights reflect the prespecified selection policy; they are not matched-coefficient comparisons.}\end{table}


All three initial checkpoints are detected on all 64 final outputs. None of the selected attacks removes a watermark, even before enforcing quality. Final preservation is 60/64 for the released model, 56/64 for the fixed-target model and 57/64 for the adapted-target model. The clean-model diagnostic is 0/64 detected. This does not establish a defense benefit: the selected attacker is ineffective against the original model too. The finite null set also cannot precisely certify a $1\%$ false-positive rate. All 384 original/attacked native PNG hashes, sample indices, labels, selection hashes and final quality hashes are checked in the analysis receipt.

\paragraph{A stronger client with explicit quality fallback.} A subsequent development protocol allows an attacker to retain its supplied marked source model at inference. It generates a source and an adapted output with coupled randomness, computes the frozen quality metrics without owner feedback, and returns the source whenever the adaptation fails quality. This hybrid client pays for both generation paths and quality evaluation. It is distinct from releasing a single adapted checkpoint. A 500-update adaptation can therefore be tested without treating its damaged samples as successful removals. Validation and final generation use fresh seeds; class identities and existing training exports are reused. This revision is frozen after the negative result above, and requires its own reporting and later confirmation.
